\chapter{Ist-Analyse}
\label{chap:ist_analyse}

\section{Beschreibung des Gesamtsystems}
\label{sec:gesamtsystem}

Die Plattform \textit{Bazario} ist ein webbasierter
Multi-Rollen-Marktplatz, der im Rahmen eines Praxisprojekts
bei SahemTech entwickelt wird. Das Gesamtsystem folgt einer
klaren Trennung zwischen Präsentationsschicht, Geschäftslogik
und Datenhaltung. Die Kommunikation zwischen Frontend und
Backend erfolgt über \acs{REST}-basierte \acs{HTTP}-Schnittstellen,
die Echtzeitkommunikation des Chat-Systems über eine
WebSocket-Verbindung auf Basis von Pusher.

Das Backend wurde von Nagham Alaatki im Rahmen einer
begleitenden Bachelorarbeit mit dem Laravel-Framework
implementiert. Es übernimmt die zentralen Funktionen der
Plattform, darunter Benutzerverwaltung, rollenbasierte
Zugriffskontrolle, Buchungsverwaltung, Werbeworkflows
sowie die Zahlungsintegration über Stripe. Die Datenhaltung
erfolgt über eine relationale MySQL-Datenbank, die
Authentifizierung über Laravel Sanctum mit tokenbasierter
\acs{API}-Authentifizierung. Die vorliegende Arbeit beschränkt
sich auf die Entwicklung des Frontends, das die vom Backend
bereitgestellten \acs{API}-Endpunkte konsumiert und alle
Funktionsbereiche der Plattform als rollenbasierte
Single-Page Application darstellt.


\section{Anforderungsanalyse}
\label{sec:anforderungsanalyse}

\subsection{Funktionale Anforderungen}
\label{subsec:funktionale-anforderungen}

Die funktionalen Anforderungen an das Frontend wurden in
enger Abstimmung mit den im Backend definierten
Geschäftsprozessen erarbeitet und gliedern sich nach den
vier Nutzerrollen der Plattform.

Alle Nutzer müssen sich registrieren, anmelden und abmelden
können. Die Registrierung erfolgt mit der Standardrolle
Customer. Nutzer können einen Antrag auf Rollenwechsel zu
Seller oder ServiceProvider stellen, indem sie ein Formular
ausfüllen und notwendige Dokumente hochladen. Der
Antragsstatus ist für den Nutzer jederzeit einsehbar.
Zusätzlich unterstützt die Plattform eine
Passwort-Zurücksetzen-Funktion auf Basis eines
\acfirst{OTP}-Verfahrens, das vom Backend bereitgestellt wird.

Käufer müssen Produkte und Dienstleistungen nach Kategorie
filtern und in einer Detailansicht einsehen können. Produkte
und Dienstleistungen können in einen gemeinsamen Warenkorb
gelegt werden. Der Checkout-Prozess umfasst die
Zahlungsabwicklung über Stripe sowie eine
Bestellbestätigungsseite. Bestellhistorie und
Buchungsübersicht müssen einsehbar und stornierbar sein.
Für Dienstleistungen müssen verfügbare Zeitfenster des
jeweiligen Anbieters angezeigt und buchbar sein.

Verkäufer benötigen ein Dashboard zur Verwaltung ihrer
Produkte einschließlich Bildupload sowie eine Übersicht
über eingegangene Bestellungen und deren Statusverwaltung.
Dienstleister müssen ihre Dienstleistungen, Arbeitszeiten
und Abwesenheiten verwalten sowie Buchungsanfragen
bestätigen oder ablehnen können. Beide Rollen können
Werbeanzeigen erstellen, die einem administrativen
Genehmigungsprozess unterliegen, und sich über
Stripe Connect für Auszahlungen registrieren.

Die administrativen Funktionen, darunter die Prüfung von
Rollenwechselanträgen, die Freigabe oder Ablehnung von
Werbeanzeigen, die Konfiguration der Plattformgebühren sowie
das Auslösen von Auszahlungen, werden über das vom Backend
bereitgestellte Admin-Dashboard abgedeckt und stellen daher
keine Anforderungen an das hier entwickelte Frontend. Alle Nutzer können Konversationen
starten und in Echtzeit Nachrichten austauschen. Die Inbox
zeigt alle aktiven Konversationen mit Ungelesen-Zähler an.

Tabelle~\ref{tab:funktionale-anforderungen} fasst die
funktionalen Anforderungen je Rolle zusammen.

\begin{table}[h]
\centering
\caption{Funktionale Anforderungen je Nutzerrolle}
\label{tab:funktionale-anforderungen}
\begin{tabular}{|p{3cm}|p{9.5cm}|}
\hline
\textbf{Rolle} & \textbf{Funktionale Anforderungen} \\
\hline
Alle Nutzer & Registrierung, Login, Logout,
Passwort zurücksetzen, Messaging, Rollenwechselantrag \\
\hline
Customer & Produktfilterung, Produktdetail,
Warenkorb, Checkout mit Stripe, Bestellhistorie,
Buchungsübersicht, Stornierung \\
\hline
Seller & Produktverwaltung (\acfirst{CRUD}), Bildupload,
Bestellübersicht, Statusverwaltung, Werbeanzeigen erstellen,
Stripe-Connect-Onboarding \\
\hline
ServiceProvider & Dienstleistungsverwaltung, Arbeitszeiten
und Abwesenheiten definieren, Buchungskalender,
Buchungsanfragen verwalten, Werbeanzeigen erstellen,
Stripe-Connect-Onboarding \\
\hline
Admin & Keine Frontend-Anforderungen; Administration über das
Backend-Dashboard (vgl. Abschnitt~\ref{sec:abgrenzung}) \\
\hline
\end{tabular}
\end{table}

\subsection{Nicht-funktionale Anforderungen}
\label{subsec:nichtfunktionale-anforderungen}

Neben den funktionalen Anforderungen wurden folgende
nicht-funktionale Anforderungen an das Frontend definiert.

Die Anwendung soll flüssig bedienbar sein. Seitenwechsel
erfolgen ohne vollständigen Seitenneuladen, Ladezeiten
werden durch Lazy Loading und effizientes State Management
minimiert. \acs{API}-Anfragen werden durch geeignete
Ladezustände und Fehlermeldungen für den Nutzer transparent
kommuniziert.

Die Sicherheit wird durch tokenbasierte Authentifizierung
über Laravel Sanctum, rollenbasierten Seitenschutz über
geschützte Routen sowie die Weiterleitung auf eine
\acs{PCI-DSS}-konforme Stripe-Zahlungsseite gewährleistet.
Sensible Daten werden zu keiner Zeit im Frontend gespeichert
oder verarbeitet. Abgelaufene Tokens werden automatisch
erkannt und der Nutzer wird zur erneuten Anmeldung
weitergeleitet.

Die Anwendung muss responsiv sein und auf gängigen Desktop-
und Mobilbrowsern korrekt dargestellt werden. Die
Benutzeroberfläche orientiert sich an den Richtlinien von
\acs{WCAG}~2.1 hinsichtlich Barrierefreiheit, Kontrast und
Tastaturnavigation. Wartbarkeit und Erweiterbarkeit werden
durch eine komponentenbasierte Architektur sowie eine
feature-orientierte Ordnerstruktur sichergestellt.

Tabelle~\ref{tab:nichtfunktionale-anforderungen} gibt
einen strukturierten Überblick über die nicht-funktionalen
Anforderungen und deren Umsetzungsmaßnahmen im Frontend.

\begin{table}[H]
\centering
\caption{Nicht-funktionale Anforderungen und Maßnahmen}
\label{tab:nichtfunktionale-anforderungen}
\begin{tabular}{|p{3cm}|p{5.3cm}|p{5cm}|}
\hline
\textbf{Anforderung} & \textbf{Beschreibung} &
\textbf{Maßnahme im Frontend} \\
\hline
Performance & Schnelle Ladezeiten, flüssige Navigation &
Lazy Loading, optimiertes State Management \\
\hline
Sicherheit & Schutz vor unbefugtem Zugriff &
Laravel Sanctum, geschützte Routen, PCI-DSS-konforme
Stripe-Zahlungsseite \\
\hline
Responsivität & Korrekte Darstellung auf allen Geräten &
Responsive Design, \acs{CSS}-Breakpoints \\
\hline
Barrierefreiheit & \acs{WCAG}~2.1-Konformität &
Semantisches \acfirst{HTML}, Kontrastverhältnisse \\
\hline
Wartbarkeit & Erweiterbare Codestruktur &
Komponentenarchitektur, feature-basierte Struktur \\
\hline
\end{tabular}
\end{table}


\section{Rollenmodell}
\label{sec:rollenmodell}

Die Plattform \textit{Bazario} implementiert ein
rollenbasiertes Zugriffskontrollmodell mit vier
Nutzerrollen. Jeder registrierte Nutzer erhält zunächst
die Rolle Customer. Ein Wechsel zu Seller oder
ServiceProvider ist durch einen Rollenwechselantrag
möglich, der durch einen Administrator genehmigt werden
muss. Der Ablauf dieses Prozesses basiert auf dem im Backend
implementierten \acs{RBAC}-Konzept. Das Frontend stellt dem
Nutzer ein Formular zum Einreichen des Antrags sowie
eine Statusanzeige bereit und passt die
Navigationselemente nach erfolgter Rollenzuweisung
dynamisch an.

Ein Nutzer kann dabei mehrere Rollen gleichzeitig
besitzen. So kann ein Nutzer beispielsweise sowohl als
Customer einkaufen als auch als Seller Produkte anbieten.
Das Frontend passt die Navigationselemente und verfügbaren
Funktionsbereiche dynamisch an die aktiven Rollen des
eingeloggten Nutzers an. Tabelle~\ref{tab:rollenmodell}
gibt einen Überblick über die Rollen und ihre jeweiligen
Funktionsbereiche im Frontend.

\begin{table}[h]
\centering
\caption{Rollenmodell der Plattform Bazario}
\label{tab:rollenmodell}
\begin{tabular}{|p{3cm}|p{9.5cm}|}
\hline
\textbf{Rolle} & \textbf{Funktionsbereiche im Frontend} \\
\hline
Customer & Produktfilterung, Warenkorb, Checkout,
Bestellhistorie, Buchungsübersicht, Messaging,
Rollenwechselantrag \\
\hline
Seller & Produktverwaltung, Bestellübersicht,
Statusverwaltung, Werbeanzeigen,
Stripe-Connect-Onboarding \\
\hline
ServiceProvider & Dienstleistungsverwaltung, Arbeitszeiten,
Abwesenheiten, Buchungskalender, Werbeanzeigen,
Stripe-Connect-Onboarding \\
\hline
Admin & Keine eigenen Frontend-Bereiche; administrative
Funktionen über das Backend-Dashboard \\
\hline
\end{tabular}
\end{table}


% =====================================================================
% ERSATZ fuer Abschnitt 3.4 in Kapitel/3_Ist_Analyse.tex
% Kompletten bisherigen Abschnitt "Technologieauswahl und Begruendung"
% durch den folgenden Text ersetzen.
%
% Ziel: Flieszttext von ca. 2 Seiten auf ca. 0,5 Seiten reduzieren.
% Die Begruendung wird von der Tabelle getragen, die Gegenueberstellung
% der Alternativen von Abschnitt 2.3.
% =====================================================================

\section{Technologieauswahl und Begründung}
\label{sec:technologieauswahl}

Die Auswahl des Frontend-Technologiestacks folgt den in
Abschnitt~\ref{sec:anforderungsanalyse} definierten Anforderungen.
Maßgeblich sind dabei vier Kriterien: die Eignung für eine
interaktionsreiche Single-Page Application mit mehreren
Nutzerrollen, die Erfüllung der nicht-funktionalen Vorgaben zu
Performance, Sicherheit, Responsivität und Barrierefreiheit
(vgl. Abschnitt~\ref{subsec:nichtfunktionale-anforderungen}), die
Kompatibilität mit dem bestehenden Laravel-Backend sowie der
Einarbeitungsaufwand im Rahmen einer Bachelorarbeit.

Die eingesetzten Technologien lassen sich anhand dieser Kriterien
in zwei Gruppen einteilen. Durch das Backend vorgegeben sind
Laravel Sanctum für die tokenbasierte Authentifizierung, Pusher
für die Echtzeitkommunikation sowie Stripe für die
Zahlungsabwicklung; für diese drei bestand keine eigenständige
Auswahlentscheidung, da das Frontend lediglich die clientseitige
Anbindung übernimmt. Im Fall von Stripe erfolgt die
Zahlungsabwicklung vollständig serverseitig: Das Frontend leitet
den Nutzer auf eine vom Backend bereitgestellte,
\acs{PCI-DSS}-konforme Zahlungsseite weiter und verarbeitet
anschließend die Rückleitung auf die Erfolgs- beziehungsweise
Abbruchseite. Eine Einbindung von Stripe Elements oder
React Stripe.js ist für den aktuellen Funktionsumfang daher nicht
erforderlich \cite{stripe_security_guide}.

Frei gewählt wurden dagegen React als Frontend-Framework, Zustand
für den clientseitigen Anwendungszustand, TanStack Query für den
serverseitigen Datenbestand, React Router für das clientseitige
Routing, Axios als \acs{HTTP}-Client, React Hook Form in Verbindung
mit Zod für Formularverarbeitung und Validierung, Tailwind CSS
mit shadcn/ui für die Gestaltung sowie Vite als Build-Tool. Die
Gegenüberstellung der jeweils in Betracht kommenden Alternativen
erfolgt in Abschnitt~\ref{sec:technologiestack}; die daraus
abgeleitete Entscheidung fasst
Tabelle~\ref{tab:technologiestack} zusammen. Für Axios war dabei
kein einzelner Grund ausschlaggebend, da die Bibliothek die im
Ökosystem am längsten bewährte Lösung für diesen Zweck darstellt
und nennenswerte Alternativen kaum vorhanden sind
\cite{axios_docs}. Die Zod-Schemata stimmen bewusst mit den
Validierungsregeln des Backends überein, sodass clientseitige und
serverseitige Prüfung denselben Regeln folgen
\cite{react_hook_form_docs, zod_docs}.

\begin{table}[H]
\centering
\caption{Technologiestack des Frontends von Bazario}
\label{tab:technologiestack}
\small
\begin{tabular}{|p{2.3cm}|p{2.4cm}|p{2.6cm}|p{4.6cm}|}
\hline
\textbf{Technologie} & \textbf{Einsatzbereich} &
\textbf{Betrachtete Alternative} & \textbf{Begründung} \\
\hline
React & \acs{UI}-Framework & Angular, Vue.js &
Komponentenbasiert, großes Ökosystem, hohe Flexibilität bei der
Strukturierung \\
\hline
Zustand & Client-State & Redux, Context \acs{API} &
Schlanker hook-basierter Store, kein Boilerplate, kein Provider
erforderlich \\
\hline
TanStack Query & Server-State & manuelle Cache-Verwaltung &
Einheitliche Verbindungszustände, Caching und gezielte
Invalidierung \\
\hline
React Router & Routing & --- &
Deklaratives Routing, rollenbasierter Seitenschutz \\
\hline
Axios & \acs{HTTP}-Kom\-mu\-ni\-ka\-tion & Fetch \acs{API} &
Zentrale Instanz mit Interceptors, automatische
Token-Anbindung \\
\hline
React Hook Form & Formular\-verwaltung & --- &
Hook-basiert, geringe Neu-Renderings, koppelbar mit
Schema-Validierung \\
\hline
Zod & Validierung & --- &
Deklarative Schemata, konsistent mit
Backend-Validierungsregeln \\
\hline
Tailwind CSS & Styling & Bootstrap, CSS-in-JS &
Utility-first, konsistentes Erscheinungsbild ohne separates
\acs{CSS} \\
\hline
shadcn/ui & \acs{UI}-Kom\-po\-nen\-ten & klassische
Komponenten\-bibliothek &
Barrierefrei, vollständig anpassbar, kein Paket-Lock-in \\
\hline
Vite & Build-Tool & --- &
Schnelle Entwicklungsumgebung, optimierte
Produktions-Builds \cite{vite_docs} \\
\hline
Laravel Sanctum & Authenti\-fizierung &
\multicolumn{2}{p{7.6cm}|}{Durch das Backend vorgegeben;
Frontend verwaltet Bearer Tokens und erkennt abgelaufene
Sitzungen \cite{laravel_sanctum_docs}} \\
\hline
pusher-js & Echtzeit\-kommunikation &
\multicolumn{2}{p{7.6cm}|}{Durch das Backend vorgegeben;
WebSocket-Anbindung über Pusher \cite{pusher_docs}} \\
\hline
Stripe (Redirect) & Zahlungs\-abwicklung &
\multicolumn{2}{p{7.6cm}|}{Durch das Backend vorgegeben;
Weiterleitung auf \acs{PCI-DSS}-konforme Zahlungsseite
\cite{stripe_security_guide}} \\
\hline
\end{tabular}
\end{table}

% =====================================================================
% TODO (Roudy) -- bitte selbst pruefen bzw. umformulieren:
%
% 1. Die vier Auswahlkriterien im ersten Absatz sind meine
%    Zusammenfassung aus Abschnitt 3.2. Sie stammen nicht aus einem
%    Meeting. Bitte pruefen, ob das deine tatsaechlichen Kriterien
%    waren, insbesondere "Einarbeitungsaufwand im Rahmen einer
%    Bachelorarbeit".
%
% 2. Spalte "Betrachtete Alternative": Folgende Eintraege sind von
%    mir ergaenzt und in deinem bisherigen Text nicht belegt:
%      - TanStack Query -> "manuelle Cache-Verwaltung"
%      - Axios -> "Fetch API"
%      - shadcn/ui -> "klassische Komponentenbibliothek"
%    Belegt sind dagegen: React -> Angular/Vue und
%    Tailwind -> Bootstrap (beide stehen so in deinem Kapitel 2).
%    Zustand -> Redux/Context ergibt sich aus dem geplanten
%    Abschnitt 2.3.3.
%    Bei React Router, React Hook Form, Zod und Vite habe ich
%    bewusst "---" gesetzt, statt Alternativen zu erfinden.
%    Jede Alternative, die hier stehen bleibt, muss in Abschnitt 2.3
%    auch tatsaechlich gegenuebergestellt werden.
%
% 3. Der Verweis \ref{sec:technologiestack} zeigt auf den bisherigen
%    Abschnitt 2.3. Nach dem Umbau von 2.3 ggf. auf die neuen
%    Unterabschnitte praezisieren.
%
% 4. \cite{vite_docs} war im alten Text vergessen und ist hier
%    ergaenzt. Ausserdem war dort "bereitstellt ." mit Leerzeichen
%    vor dem Punkt.
% =====================================================================